\chapter{文献标注}
论文正文的引用标注应遵照原模板所列GB/T 7714。原模板允许顺序编码制或著者出版年制，但全文须统一。本项目选用顺序编码制，参考文献按首次引用顺序自动排列。
\section{引用单篇文献}
图书引用示例：数学思想可用于理解信息技术中的基础问题\cite{wu2012}。期刊论文引用示例：三维重建技术的应用可参考相关综述\cite{ma2018}。
\section{引用多篇文献}
在同一处引用多篇文献时，将文献键写入同一条引用命令\cite{wu2012,ma2018,chen2022}。连续编号按引用样式自动压缩，避免手动输入序号。
\section{参考文献数据维护}
在参考文献数据库中填写真实作者、题名、出版信息和页码。英文人名采用姓在前、名的首字母在后的形式；英文论文题名按句首大写处理，专有名词用花括号保护。期刊名称的大小写应在数据库中按原模板要求录入。

对于数据集、预印本、档案等特殊类型，可使用数据库的mark字段指定DS、PP、A等标识，medium字段指定OL等载体标识；具体字段示例见项目说明。正式提交前仍需逐条核对参考文献数据与实际引用。
\section{本章小结}
正文只列出实际引用的文献。未被引用但被阅读或具有补充信息的文献，可按原模板要求列入附录中的“参考书目”。
